% Auto-converted from paper2_full_draft.md §2.
% Wording preserved; no claims added.
\section{Introduction}
Vulnerability triage at scale relies on signals that are easy to share: the \textbf{Exploit Prediction Scoring System (EPSS)}~\cite{Jacobs2021EPSS,Jacobs2023EnhancingEPSS}, the \textbf{CISA Known Exploited Vulnerabilities (KEV) catalog}, \textbf{CVSS v3.1 / v4} severity, and local context such as asset criticality, network exposure, and remediation cost. Practitioners frequently combine these signals with hand-tuned or fitted weights, and companion work explores augmenting EPSS-weighted prioritization with cyber-hygiene signals on related slices~\cite{paper4hygieneprio}. The fitted-weight path is attractive because calibration promises an estimator of risk grounded in observed exploitation outcomes. It is also potentially fragile: per-feature weight estimation requires \emph{enough labelled positives}, not enough labelled \emph{pairs} (the calibration unit is the CVE, not the host-CVE pair), and the available public exploit signals are sparse on time-bounded windows.

This paper is a companion to the VulnPrio study, which contributed a benchmark and scheduler for context-aware prioritization on a synthetic fleet with placeholder weights; a separate companion contributes a reproducible synthetic benchmark for the hygiene-anomaly setting~\cite{paper3hygienebench}. The VulnPrio study made no calibration or model-quality claim. The present paper asks whether public-feed signals can move us from placeholder weights toward weights with a calibrated interpretation, a question that companion work also pursues from the angle of rolling-history online calibration~\cite{paper7online}. The honest answer is no, at the slice we examined. The negative finding is not a bug; it is a data-density boundary that needs to be acknowledged before public-feed calibration becomes a default move. The unit of analysis for per-feature weight learning is the distinct CVE-level positive, not the host-CVE record.

The central contribution of this paper, stated crisply, is a \emph{failure-aware gating mechanism} for public-feed vulnerability prioritization: a pre-registered, reproducible procedure that detects when an exploit-likelihood signal's calibration has failed for lack of distinct labelled positives and, on that detection, falls back to a safe fixed-prior ranking instead of emitting an unsupported calibrated score. The operational payoff is that the prioritizer never silently substitutes a mis-calibrated estimator for an honest one under the sparse-label conditions that public-sector feeds routinely impose. What this paper adds beyond prior public-feed prioritization work is twofold: an explicit, measured calibration-feasibility boundary on a realistic catalog-restricted slice, the distinct-positive density at which the calibration framing ceases to be defensible, and the gate that enforces the safe fallback once that boundary is crossed. We then describe the gate, the sensitivity framework, the measured gate outcome, and the confirmatory sensitivity results as a self-contained study. We position the work against EPSS / KEV / CVSS / SSVC and against recent capacity-constrained prioritization, EPSS-stability, and context-feature-ablation studies. We do not claim that any fixed prior matches or exceeds EPSS in any general sense. We do not claim deployment readiness. We do claim that the gate methodology is reproducible and that the negative feasibility result is informative.
